\section{Conclusion}
\label{sec:conclusion}

The legal case for algorithmic redistricting is strong on every dimension examined in this Article. The constitutional authority exists---the Elections Clause, as construed in \textit{Smiley v.\ Holm}, permits Congress to mandate the method of congressional redistricting, and the Huntington-Hill statute provides a direct structural precedent for mathematical mandates in the apportionment-redistricting domain. The state adoption pathways are clear---42 states can adopt algorithmic redistricting by ordinary statute or commission resolution without constitutional amendment, and the remaining eight face modest amendment requirements. The VRA compliance case is strong---the algorithm produces more majority-minority districts than enacted plans, and its VRA mode creates majority-minority districts through documented compactness optimization rather than racial targeting, satisfying the \textit{Gingles} framework while avoiding the \textit{Shaw} racial gerrymandering concern. And the anticipated legal challenges do not hold up: the nondelegation claim mischaracterizes delegation to an executive agency implementing a statutory mandate as delegation to the algorithm itself; the Elections Clause limits claim was resolved in 1932; and the VRA conflict claim gets the empirical facts backwards.

The model statute in Section~\ref{sec:model} provides ready-to-use text for federal or state enactment. Its three-section structure---definitions and certification, map production and adjustment, judicial review---can be adapted to any state's existing redistricting statutory scheme with minimal modification. State legislatures working from this template should be attentive to two adaptation requirements: substituting the appropriate state agency for the Bureau of the Census, and ensuring that the adjustment criteria are consistent with any constitutional redistricting standards their state constitution specifies.

The post-\textit{Rucho} moment is, paradoxically, an opportunity. The federal courthouse is closed to partisan gerrymandering claims. State courts are active but inconsistent. Congress has an explicit invitation from the Chief Justice to act. The tool is ready, the legal authority is clear, and the public supports reform.\footnote{See \citet{dellaLibera2026survey} (finding 71\% of survey respondents support their state using computer algorithms to draw congressional districts).} The question is whether the political will exists to use them.

The experience of the Huntington-Hill apportionment method should be instructive. Before 1941, apportionment was among the most contentious recurring political disputes in American governance. After 1941, the mathematical formula rendered the question technical, predictable, and uncontroversial. Eighty years of operational success is strong evidence that mathematical mandates can resolve politically charged democratic governance questions---not by choosing winners and losers, but by ensuring that the process of choosing cannot itself be weaponized. Congressional redistricting is overdue for the same resolution.

\bigskip

\noindent\textit{The author thanks the algorithmic redistricting research team for technical collaboration and data access. All errors are the author's own.}
